\documentclass[12pt,a4paper]{article}
\usepackage[UTF8]{ctex}
\usepackage{amsmath,amssymb,amsthm}
\usepackage{geometry}
\usepackage{bm}
\usepackage{hyperref}
\usepackage{graphicx}

\geometry{margin=2.5cm}

\title{第12章\quad 抓取与操作}
\author{现代机器人学：力学、规划与控制\\
Kevin M. Lynch 和 Frank C. Park\\
中文翻译}
\date{}

\begin{document}

\maketitle

\tableofcontents
\newpage

\section{引言}

前面章节的重点是描述、规划和控制机器人本身的运动。为了完成有用的工作，机器人必须能够操作环境中的物体。在本章中，我们建模机器人与物体之间的接触，特别是接触对物体运动施加的约束以及可以通过摩擦接触传递的力。利用这些模型，我们研究通过形状闭合（form closure）和力闭合（force closure）抓取来选择接触以固定物体的问题。我们还将接触建模应用于抓取以外的操作问题，例如推动物体、动态携带物体以及测试结构的稳定性。

\section{接触运动学}
\label{sec:contact_kinematics}

接触运动学是研究两个或多个刚体在遵守不可穿透约束的情况下如何相对运动。它还将接触处的运动分类为滚动或滑动。让我们从研究两个刚体之间的单个接触开始。

\subsection{单个接触的一阶分析}
\label{subsec:first_order_analysis}

考虑两个刚体，其配置分别由局部坐标列向量$q_1$和$q_2$给出。将复合配置写为$q = (q_1, q_2)$，我们定义物体之间的距离函数$d(q)$，当它们分离时为正，接触时为零，穿透时为负。当$d(q) > 0$时，物体的运动没有约束；每个物体都可以自由地以六个自由度运动。当物体接触时（$d(q) = 0$），我们查看时间导数$\dot{d}$、$\ddot{d}$等，以确定物体在遵循特定轨迹$q(t)$时是保持接触还是分离。这可以通过以下可能性表确定：

\begin{center}
\begin{tabular}{|c|c|c|c|}
\hline
$d$ & $\dot{d}$ & $\ddot{d}$ & $\cdots$ \\
\hline
$>0$ & & & 无接触 \\
$<0$ & & & 不可行（穿透） \\
$=0$ & $>0$ & & 接触，但正在分离 \\
$=0$ & $<0$ & & 不可行（穿透） \\
$=0$ & $=0$ & $>0$ & 接触，但正在分离 \\
$=0$ & $=0$ & $<0$ & 不可行（穿透） \\
\hline
\end{tabular}
\end{center}

只有当所有时间导数都为零时，接触才得以保持。

现在假设两个物体最初在单点接触（$d = 0$）。$d$的前两个时间导数写为
\begin{equation}
\dot{d} = \frac{\partial d}{\partial q} \dot{q},
\label{eq:12.1}
\end{equation}
\begin{equation}
\ddot{d} = \dot{q}^T \frac{\partial^2 d}{\partial q^2} \dot{q} + \frac{\partial d}{\partial q} \ddot{q}.
\label{eq:12.2}
\end{equation}

项$\partial d/\partial q$和$\partial^2 d/\partial q^2$携带关于局部接触几何的信息。梯度向量$\partial d/\partial q$对应于与接触法线相关的$q$空间中的分离方向（图12.2）。矩阵$\partial^2 d/\partial q^2$编码接触点处物体的相对曲率信息。

在本章中，我们假设在接触处只有接触法线信息$\partial d/\partial q$可用；关于局部接触几何的其他信息，包括接触曲率$\partial^2 d/\partial q^2$和更高阶导数，是未知的。基于这个假设，我们在方程(12.1)处截断分析，并假设如果$\dot{d} = 0$，物体保持接触。由于我们只处理一阶接触导数$\partial d/\partial q$，我们将分析称为一阶分析。在这种一阶分析中，图12.2中的接触被视为相同，因为它们具有相同的接触法线。

如上面的表所示，包含接触曲率$\partial^2 d/\partial q^2$的二阶分析可能表明，即使$d = \dot{d} = 0$，接触实际上也在分离或穿透。我们将看到这样的例子，但二阶接触条件的详细分析超出了本章的范围。

\subsection{接触类型：滚动、滑动和分离}
\label{subsec:contact_types}

给定两个单点接触的物体，如果接触得以保持，它们会经历滚动-滑动运动。保持接触的约束是完整约束，$d(q) = 0$。保持接触的必要条件是$\dot{d} = 0$。

让我们基于接触法线，以不显式需要距离函数的形式写出速度约束$\dot{d} = 0$（图12.2）。设$\hat{n} \in \mathbb{R}^3$是与接触法线对齐的单位向量，在世界坐标系中表示。设$p_A \in \mathbb{R}^3$是物体$A$上接触点在世界坐标系中的表示，设$p_B \in \mathbb{R}^3$是物体$B$上接触点在世界坐标系中的表示。虽然接触点向量$p_A$和$p_B$最初是相同的，但速度$\dot{p}_A$和$\dot{p}_B$可能不同。因此，条件$\dot{d} = 0$可以写为
\begin{equation}
\hat{n}^T (\dot{p}_A - \dot{p}_B) = 0.
\label{eq:12.3}
\end{equation}

由于接触法线的方向定义为指向物体$A$，不可穿透约束$\dot{d} \geq 0$写为
\begin{equation}
\hat{n}^T (\dot{p}_A - \dot{p}_B) \geq 0.
\label{eq:12.4}
\end{equation}

让我们用物体$A$和$B$在空间坐标系中的twist $V_A = (\omega_A, v_A)$和$V_B = (\omega_B, v_B)$重写约束(12.4)。\footnote{在本章中，所有twist和wrench都在空间坐标系中表示。}注意
\begin{align}
\dot{p}_A &= v_A + \omega_A \times p_A = v_A + [\omega_A]p_A, \\
\dot{p}_B &= v_B + \omega_B \times p_B = v_B + [\omega_B]p_B.
\end{align}

我们还可以定义对应于沿接触法线施加的单位力的wrench $F = (m, f)$：
\begin{equation}
F = (p_A \times \hat{n}, \hat{n}) = ([p_A]\hat{n}, \hat{n}).
\end{equation}

在刚体的纯运动学分析中，没有必要诉诸力，但我们现在采用这种记号是方便的，以预期第12.2节中关于接触力的讨论。

利用这些表达式，不等式约束(12.4)可以写为
\begin{equation}
\text{（不可穿透约束）} \quad F^T (V_A - V_B) \geq 0.
\label{eq:12.5}
\end{equation}

如果
\begin{equation}
\text{（主动约束）} \quad F^T (V_A - V_B) = 0,
\label{eq:12.6}
\end{equation}
那么，在一阶情况下，约束是主动的，物体保持接触。

在$B$是静止夹具的情况下，不可穿透约束(12.5)简化为
\begin{equation}
F^T V_A \geq 0.
\label{eq:12.7}
\end{equation}

如果$F^T V_A > 0$，则称$F$和$V_A$是排斥的。如果$F^T V_A = 0$，则称$F$和$V_A$是互反的，约束是主动的。

满足(12.6)的twist $V_A$和$V_B$称为一阶滚动-滑动运动——接触可能是滑动或滚动。滚动-滑动接触可以进一步分为滚动接触和滑动接触。如果物体在接触处没有相对运动，则接触是滚动的：
\begin{equation}
\text{（滚动约束）} \quad \dot{p}_A = v_A + [\omega_A]p_A = v_B + [\omega_B]p_B = \dot{p}_B.
\label{eq:12.8}
\end{equation}

注意"滚动"接触包括两个物体保持相对静止的情况，即没有相对旋转。因此，"粘附"是这些接触的另一个术语。

如果twist满足方程(12.6)但不满足(12.8)的滚动方程，则它们是滑动的。

我们为滚动接触分配接触标签$R$，为滑动接触分配标签$S$，为正在分离的接触（满足不可穿透约束(12.5)但不满足主动约束(12.6)）分配标签$B$。

滚动和滑动接触之间的区别在我们考虑第12.2节中的摩擦力时变得特别重要。

\subsection{多个接触}
\label{subsec:multiple_contacts}

现在假设物体$A$受到与$m$个其他物体的$n$个接触，其中$n \geq m$。接触编号为$i = 1, \ldots, n$，其他物体编号为$j = 1, \ldots, m$。设$j(i) \in \{1, \ldots, m\}$表示参与接触$i$的其他物体的编号。每个接触$i$将$V_A$约束到其六维twist空间的半空间，该半空间由形式为$F_i^T V_A = F_i^T V_{j(i)}$的五维超平面界定。取所有接触的约束集的并集，我们在$V_A$空间中得到可行twist的多面体凸集（简称多面体）$\mathcal{V}$，写为
\begin{equation}
\mathcal{V} = \{V_A | F_i^T (V_A - V_{j(i)}) \geq 0 \text{ 对所有 } i\},
\end{equation}
其中$F_i$对应于第$i$个接触法线（指向物体$A$），$V_{j(i)}$是接触$i$处其他物体的twist。如果接触$i$贡献的半空间约束不改变可行twist多面体$\mathcal{V}$，则接触$i$处的约束是冗余的。通常，物体的可行twist多面体可以包括六维内部（没有接触约束是主动的）、一个约束主动的五维面、两个约束主动的四维面，依此类推，直到一维边和零维点。在多面体的$k$维面上的twist $V_A$表明有$6 - k$个独立（非冗余）的接触约束是主动的。

如果提供约束的所有物体都是静止的，即对所有$j$有$V_j = 0$，则由(12.5)定义的每个约束超平面都通过$V_A$空间的原点。我们称这种约束为齐次的。可行twist集成为以原点为根部的锥，称为（齐次）多面体凸锥。设$F_i$是静止接触$i$的约束wrench。那么可行twist锥$\mathcal{V}$是
\begin{equation}
\mathcal{V} = \{V_A | F_i^T V_A \geq 0 \text{ 对所有 } i\}.
\end{equation}

如果$F_i$正张成六维wrench空间，或者等价地，$F_i$的凸包包含原点在其内部，则可行twist多面体$\mathcal{V}$简化为原点处的点，静止接触完全约束物体的运动，我们得到形状闭合，在第12.1.7节中更详细地讨论。

如第12.1.2节所述，每个点接触$i$可以给定一个对应于接触类型的标签：如果接触正在分离，则为$B$；如果接触是滚动的，则为$R$；如果接触是滑动的，即满足(12.6)但不满足(12.8)，则为$S$。整个系统的接触模式可以写为接触处接触标签的连接。由于我们有三个不同的接触标签，具有$n$个接触的物体系统最多可以有$3^n$个接触模式。其中一些接触模式可能不可行，因为它们对应的运动学约束可能不兼容。

\subsection{物体集合}
\label{subsec:collections_of_bodies}

上面的讨论可以推广到寻找多个接触物体的可行twist。如果物体$i$和$j$在点$p$处接触，其中$\hat{n}$指向物体$i$且$F = ([p]\hat{n}, \hat{n})$，则它们的空间twist $V_i$和$V_j$必须满足约束
\begin{equation}
F^T (V_i - V_j) \geq 0
\label{eq:12.13}
\end{equation}
以避免穿透。这是复合$(V_i, V_j)$ twis t空间中的齐次半空间约束。在$m$个物体的装配中，每个成对接触在$6m$维复合twist空间（对于平面物体为$3m$维）中贡献另一个约束，结果是运动学可行twist的以复合twist空间原点为根部的多面体凸锥。整个装配的接触模式是装配中每个接触处接触标签的连接。

如果有运动受控的物体，例如机器人手指，则对剩余物体运动的约束不再是齐次的。因此，未控制物体在其复合twist空间中的凸多面体可行twist集不再是以原点为根部的锥。

\subsection{其他类型的接触}
\label{subsec:other_types_of_contacts}

我们一直在考虑图12.6(a)所示的点接触类型，其中接触的至少一个物体唯一地定义了接触法线。图12.6(b)-(e)显示了其他类型的接触。图12.6(b)-(d)的凸-凹顶点、线和平面接触提供的运动学约束，在一阶情况下，与有限个单点接触集合提供的约束相同。图12.6(e)中的退化情况被忽略，因为没有接触法线的唯一定义。

不可穿透约束(12.5)源于可以在法线方向施加任意大的接触力以防止穿透的事实。在第12.2节中，我们将看到由于摩擦而产生的切向力也可以施加，这些力可以防止两个接触物体之间的滑动。法向和切向接触力受到约束：法向力必须推入物体，而不是拉出，最大摩擦力与法向力成正比。

如果我们希望应用可以近似摩擦效果而不显式建模力的运动学分析，我们可以定义三种纯运动学点接触模型：无摩擦点接触、有摩擦点接触和软接触，也称为软指接触。无摩擦点接触只强制执行滚动-滑动约束(12.5)。有摩擦点接触还强制执行滚动约束(12.8)，隐式建模足以防止接触处滑动的摩擦力。软接触强制执行滚动约束(12.8)以及一个额外的约束：两个接触物体可能不会绕接触法线轴相对旋转。这模拟了变形和由于两个物体之间的非零接触面积而产生的抵抗任何旋转的摩擦力矩。对于平面问题，有摩擦点接触和软接触是相同的。

\subsection{平面图形方法}
\label{subsec:planar_graphical_methods}

平面问题允许使用图形方法来可视化单个物体的可行运动，因为twist空间是三维的。图12.5显示了一个平面twist锥的例子。对于具有超过三个自由度的系统，这样的图将非常难以绘制。

表示平面twist $V = (\omega_z, v_x, v_y)$在$\{s\}$中的一种方便方法是作为旋转中心（CoR）在$(-v_y/\omega_z, v_x/\omega_z)$加上角速度$\omega_z$。CoR是在运动下保持静止的（投影）平面中的点，即螺旋轴与平面相交的点。\footnote{注意，$\omega_z = 0$的情况必须小心处理，因为它对应于无穷远处的CoR。}在运动速度不重要的情况下，我们可以简单地用表示旋转方向的$'+'$、$'-'$或$0$符号标记CoR（图12.7）。从平面twist到CoR的映射如图12.8所示，它显示CoR空间由$'+'$ CoR（逆时针）的平面、$'-'$ CoR（顺时针）的平面和平移方向的圆组成。

给定两个不同的twist $V_1$和$V_2$及其对应的CoR，这些twist的线性组合集，$k_1 V_1 + k_2 V_2$，其中$k_1, k_2 \in \mathbb{R}$，对应于通过$\text{CoR}(V_1)$和$\text{CoR}(V_2)$的CoR线。由于$k_1$和$k_2$可以有任一符号，如果$\omega_{1z}$或$\omega_{2z}$非零，则此线上的CoR可以有任一符号。如果$\omega_{1z} = \omega_{2z} = 0$，则线性组合集对应于所有平移方向的集。

更有趣的情况是当$k_1, k_2 \geq 0$时。给定两个twist $V_1$和$V_2$，这两个速度的非负线性组合写为
\begin{equation}
V = \text{pos}(\{V_1, V_2\}) = \{k_1 V_1 + k_2 V_2 | k_1, k_2 \geq 0\},
\end{equation}
这是以原点为根部的平面twist锥，$V_1$和$V_2$定义锥的边。如果$\omega_{1z}$和$\omega_{2z}$具有相同的符号，则它们的非负线性组合的CoR $\text{CoR}(\text{pos}(\{V_1, V_2\}))$都具有该符号并位于两个CoR之间的线段上。如果$\text{CoR}(V_1)$和$\text{CoR}(V_2)$分别标记为$'+'$和$'-'$，则$\text{CoR}(\text{pos}(\{V_1, V_2\}))$由包含两个CoR的线组成，减去CoR之间的线段。这个集由连接到$\text{CoR}(V_1)$的标记为$'+'$的CoR射线、连接到$\text{CoR}(V_2)$的标记为$'-'$的CoR射线和对应于平移的标记为$0$的无穷远点组成。这个元素集合应该被视为单个线段（尽管一个通过无穷远），就像上面提到的第一种情况一样。图12.9和12.10显示了对应于平面twist的非负线性组合的CoR区域的例子。

平面twist的CoR表示对于表示与静止物体接触的一个可动物体的可行运动特别有用。由于约束是静止的，如第12.1.3节所述，可行twist形成以原点为根部的多面体凸锥。这样的锥可以唯一地由具有$'+'$、$'-'$和$0$标签的CoR集表示。由移动约束生成的一般twist多面体不能由具有此类标签的CoR集唯一表示。

给定静止物体和可动物体之间的接触，我们可以绘制不违反不可穿透约束的CoR。标记接触法线上所有点为$\pm$，内法线左侧的点为$'+'$，右侧的点为$'-'$。所有标记为$'+'$的点可以作为可动物体正角速度的CoR，所有标记为$'-'$的点可以作为负角速度的CoR，而不违反一阶接触约束。我们可以进一步为每个CoR分配接触标签，对应于分离接触$B$、滑动接触$S$或滚动接触$R$的一阶条件。对于平面滑动，我们将标签$S$细分为两个子类：$S_r$，其中可动物体相对于固定约束向右滑动；$S_l$，其中可动物体相对于固定约束向左滑动。图12.11说明了标记。

如果有多个接触，我们只需取各个接触的约束和接触标签的并集。约束的这种并集意味着可行CoR区域是凸的，就像齐次多面体twist锥一样。

\subsection{形状闭合}
\label{subsec:form_closure}

如果一组静止约束阻止物体的所有运动，则实现物体的形状闭合。如果这些约束由机器人手指提供，我们称之为形状闭合抓取。图12.13显示了一个例子。

\subsubsection{一阶形状闭合所需的接触数量}
\label{subsubsec:number_of_contacts_form_closure}

每个静止接触$i$提供形式为
\begin{equation}
F_i^T V \geq 0
\end{equation}
的半空间twist约束。

如果满足约束的唯一twist $V$是零twist，则形状闭合成立。对于$j$个接触，这个条件等价于
\begin{equation}
\text{pos}(\{F_1, \ldots, F_j\}) = \mathbb{R}^6
\end{equation}
对于三维物体。因此，根据本章开头的性质2，至少需要$6 + 1 = 7$个接触才能实现三维物体的一阶形状闭合。对于平面物体，条件是
\begin{equation}
\text{pos}(\{F_1, \ldots, F_j\}) = \mathbb{R}^3,
\end{equation}
需要$3 + 1 = 4$个接触才能实现一阶形状闭合。这些结果总结在以下定理中。

\begin{theorem}[12.6]
对于平面物体，至少需要四个点接触才能实现一阶形状闭合。对于三维物体，至少需要七个点接触。
\end{theorem}

现在考虑在平面中抓取圆盘的问题。应该清楚的是，无论接触数量如何，运动学上阻止圆盘的运动是不可能的；它总是能够绕其中心旋转。这样的物体被称为异常的——物体上所有点的接触法向力的正张成不等于$\mathbb{R}^n$，其中$n = 3$在平面情况下，$n = 6$在三维情况下。三维中此类物体的例子包括旋转表面，如球体和椭球体。

图12.14显示了平面抓取的例子。第12.1.6节的图形方法表明，图12.14(a)中的四个接触固定了物体。我们的一阶分析表明，图12.14(b)和12.14(c)中的物体在三指抓取中都可以绕其中心旋转，但实际上这对于图12.14(b)中的物体是不可能的——二阶分析会告诉我们这个物体实际上是被固定的。最后，一阶分析告诉我们，图12.14(d)-(f)中的两指抓取是相同的，但实际上图12.14(f)中的物体仅由两个手指固定，这是由于曲率效应。

总结一下，我们的一阶分析总是正确地标记分离和穿透运动，但二阶和更高阶效应可能将一阶滚动-滑动运动改变为分离或穿透。如果物体通过一阶分析处于形状闭合，则它在任何分析中都处于形状闭合。如果只有滚动-滑动运动通过一阶分析是可行的，则物体可能通过更高阶分析处于形状闭合；否则，物体在任何分析中都不处于形状闭合。

\subsubsection{一阶形状闭合的线性规划测试}
\label{subsubsec:linear_programming_test}

设$F = [F_1 \, F_2 \, \cdots \, F_j] \in \mathbb{R}^{n \times j}$是一个矩阵，其列由$j$个接触wrench形成。对于三维物体，$n = 6$；对于平面物体，$n = 3$，$F_i = [m_{iz} \, f_{ix} \, f_{iy}]^T$。如果存在权重向量$k \in \mathbb{R}^j$，$k \geq 0$，使得对所有$F_{\text{ext}} \in \mathbb{R}^n$有$F k + F_{\text{ext}} = 0$，则接触产生形状闭合。

显然，如果$F$的秩不是满的（$\text{rank}(F) < n$），则物体不处于形状闭合。如果$F$是满秩的，则形状闭合条件等价于存在严格正系数$k > 0$使得$F k = 0$。我们可以将这个测试表述为以下条件集，这是线性规划的一个例子：
\begin{align}
\text{寻找 } k & \\
\text{最小化 } \mathbf{1}^T k & \\
\text{使得 } F k = 0 & \\
k_i \geq 1, \quad i = 1, \ldots, j, &
\end{align}
其中$\mathbf{1}$是$j$维全1向量。如果$F$是满秩的且存在(12.14)的解$k$，则物体处于一阶形状闭合。否则不是。注意，目标函数$\mathbf{1}^T k$对于回答二元问题不是必需的，取决于LP求解器，但包含它以确保问题适定。

\subsubsection{测量形状闭合抓取的质量}
\label{subsubsec:measuring_quality_form_closure}

考虑图12.16中显示的两个形状闭合抓取。哪个是更好的抓取？

回答这个问题需要一个度量来测量抓取的质量。抓取度量取接触集$\{F_i\}$并返回单个值$\text{Qual}(\{F_i\})$，其中$\text{Qual}(\{F_i\}) < 0$表示抓取不产生形状闭合，较大的正值表示更好的抓取。

有许多合理的抓取度量选择。作为一个例子，假设为了避免损坏物体，我们要求接触$i$处的力的大小小于或等于$f_{i,\max} > 0$。那么$j$个接触可以施加的接触wrench的总集由下式给出：
\begin{equation}
\mathcal{C}_F = \left\{\sum_{i=1}^j f_i F_i \middle| f_i \in [0, f_{i,\max}] \right\}.
\label{eq:12.15}
\end{equation}

图12.17显示了二维中的例子。这显示了接触可以施加的wrench的凸集，以抵抗施加到物体的干扰wrench。如果抓取产生形状闭合，则该集在其内部包含wrench空间的原点。

现在的问题是将这个多面体转换为表示抓取质量的单个数字。理想情况下，这个过程会使用物体可能遇到的干扰wrench的某种想法。更简单的选择是将$\text{Qual}(\{F_i\})$设置为以wrench空间原点为中心、适合凸多面体内部的最大wrench球的半径。在评估这个半径时，应该考虑两个注意事项：(1)力矩和力具有不同的单位，因此没有明显的方法来等同力和力矩的大小，以及(2)由于接触力而产生的力矩取决于空间坐标系原点的位置。为了解决(1)，通常选择被抓取物体的特征长度$r$并将接触力矩$m$转换为力$m/r$。为了解决(2)，原点可以选择在物体几何中心附近或在其质心处。

给定空间坐标系和特征长度$r$的选择，我们简单地计算从wrench空间原点到$\mathcal{C}_F$边界上每个超平面的有符号距离。这些距离的最小值是$\text{Qual}(\{F_i\})$（图12.17）。

回到我们在图12.16中的原始例子，我们可以看到，如果允许每个手指施加相同的力，则左侧的抓取可能被认为更好，因为接触可以抵抗关于物体中心的更大力矩。

\subsubsection{选择形状闭合的接触}
\label{subsubsec:choosing_contacts_form_closure}

已经提出了许多方法来选择用于固定或抓取的形状闭合接触。一种方法是在物体表面上采样候选抓取点（平面物体为四个，三维物体为七个），直到找到产生形状闭合的集合。从那里，可以根据梯度上升增量地重新定位候选抓取点，使用抓取度量，即$\partial \text{Qual}(p)/\partial p$，其中$p$是所有接触位置坐标的向量。

\section{接触力和摩擦}
\label{sec:contact_forces_friction}

\subsection{摩擦}
\label{subsec:friction}

机器人操作中常用的摩擦模型是库仑摩擦。这个实验定律指出，切向摩擦力大小$f_t$与法向力大小$f_n$相关，关系为$f_t \leq \mu f_n$，其中$\mu$称为摩擦系数。如果接触正在滑动，或当前正在滚动但即将滑动（即，在下一瞬间接触正在滑动），则$f_t = \mu f_n$，摩擦力的方向与滑动方向相反，即摩擦耗散能量。摩擦力与滑动速度无关。

通常定义两个摩擦系数，静摩擦系数$\mu_s$和动摩擦（或滑动）摩擦系数$\mu_k$，其中$\mu_s \geq \mu_k$。这意味着有更大的摩擦力可用于抵抗初始运动，但一旦运动开始，阻力更小。已经开发了许多其他摩擦模型，对诸如滑动速度和滑动前静态接触持续时间等因素具有不同的函数依赖性。这些都是复杂微观行为的聚合模型。为简单起见，我们将使用具有单个摩擦系数$\mu$的最简单的库仑摩擦模型。这个模型对于硬、干燥的材料是合理的。摩擦系数取决于接触的两种材料，通常在0.1到1之间。

对于法线在$+\hat{z}$方向的接触，可以通过接触传递的力集满足
\begin{equation}
\sqrt{f_x^2 + f_y^2} \leq \mu f_z, \quad f_z \geq 0.
\label{eq:12.16}
\end{equation}

图12.18(a)显示这组力形成摩擦锥。手指可以施加到平面的力集位于所示的锥内。图12.18(b)从侧视图显示相同的锥，说明摩擦角$\alpha = \tan^{-1} \mu$，这是锥的半角。如果接触不滑动，力可以在锥内的任何地方。如果手指向右滑动，它施加的力位于摩擦锥的右边缘，大小由法向力确定。相应地，平面对手指施加相反的力，该力的切向（摩擦）部分的方向与滑动方向相反。

为了允许接触力学问题的线性表述，通常方便地用多面体凸锥表示凸圆锥。图12.18(c)显示了摩擦锥的内接四边锥形近似，由锥边的正张成定义：$(\mu, 0, 1)$、$(-\mu, 0, 1)$、$(0, \mu, 1)$和$(0, -\mu, 1)$。我们可以通过使用更多边来获得对圆锥的更紧密近似。内接锥低估了可用的摩擦力，而外接锥高估了摩擦力。选择使用哪一个取决于应用。例如，如果我们想确保机器人手可以抓取物体，低估可用的摩擦力是一个好主意。

对于平面问题，不需要近似——摩擦锥由锥的两条边的正张成精确表示，类似于图12.18(b)所示的侧视图。

一旦我们选择坐标系，任何接触力都可以表示为wrench $F = ([p]f, f)$，其中$p$是接触位置。这将摩擦锥转换为wrench锥。图12.19显示了一个平面例子。平面摩擦锥的两条边在wrench空间中给出两条射线，可以通过接触传递到物体的wrench给出沿这些边的基向量的正张成。如果$F_1$和$F_2$是这些wrench锥边的基向量，我们将wrench锥写为$\mathcal{WC} = \text{pos}(\{F_1, F_2\})$。

如果多个接触作用在物体上，则可以通过接触传递到物体的总wrench集是所有单个wrench锥$\mathcal{WC}_i$的正张成：
\begin{equation}
\mathcal{WC} = \text{pos}(\{\mathcal{WC}_i\}) = \left\{\sum_i k_i F_i \middle| F_i \in \mathcal{WC}_i, k_i \geq 0 \right\}.
\end{equation}

这个复合wrench锥是以原点为根部的凸锥。图12.19(d)显示了具有图12.19(c)所示两个摩擦锥的平面物体的这种复合wrench锥的例子。对于平面问题，三维wrench空间中的复合wrench锥是多面体的。对于三维问题，六维wrench空间中的wrench锥不是多面体的，除非单个摩擦锥由多面体锥近似，如图12.18(c)所示。

如果作用在物体上的接触或接触集是理想力控制的，则控制器指定的wrench $F_{\text{cont}}$必须位于对应于这些接触的复合wrench锥内。如果还有其他非力控制的接触作用在物体上，则物体上可能的wrench锥等价于来自非力控制接触的wrench锥，但平移到以$F_{\text{cont}}$为根部。

\subsection{平面图形方法}
\label{subsec:planar_graphical_methods_force}

\subsubsection{表示Wrench}
\label{subsubsec:representing_wrenches}

任何具有非零线性分量的平面wrench $F = (m_z, f_x, f_y)$可以表示为平面中绘制的箭头，其中箭头的基座在点
\begin{equation}
(x, y) = \frac{1}{f_x^2 + f_y^2} (m_z f_y, -m_z f_x)
\end{equation}
箭头的头部在$(x + f_x, y + f_y)$。如果我们在其线上的任何地方滑动箭头，力矩不变，因此沿此线的相同方向和长度的任何箭头表示相同的wrench（图12.20）。如果$f_x = f_y = 0$且$m_z \neq 0$，则wrench是纯力矩，我们不尝试用图形表示它。

两个表示为箭头的wrench可以通过沿其线滑动箭头直到箭头的基座重合来图形求和。两个wrench之和对应的箭头如图12.20所示。该方法可以顺序应用于对表示为箭头的多个wrench求和。

\subsubsection{表示Wrench锥}
\label{subsubsec:representing_wrench_cones}

在上一节中，每个wrench都有指定的大小。然而，刚体接触意味着接触法向力可以任意大；法向力达到防止两个物体穿透所需的大小。因此，拥有所有形式为$kF$的wrench的表示是有用的，其中$k \geq 0$且$F \in \mathbb{R}^3$是基向量。

一种这样的表示是力矩标记。基wrench $F$的箭头按照第12.2.2.1节所述绘制。然后，箭头线左侧平面中的所有点标记为$'+'$，表示$F$的任何正缩放对这些点产生正力矩$m_z$，箭头线右侧平面中的所有点标记为$'-'$，表示$F$的任何正缩放对这些点产生负力矩。线上的点标记为$\pm$。

推广，力矩标记可以表示任何齐次凸平面wrench锥，就像齐次凸平面twist锥可以表示为凸CoR区域一样。给定对应于所有$k_i \geq 0$的wrench $k_i F_i$的定向力线集合，wrench锥$\text{pos}(\{F_i\})$可以通过标记平面中的每个点来表示：如果每个$F_i$对该点产生非负力矩，则标记为$'+'$；如果每个$F_i$对该点产生非正力矩，则标记为$'-'$；如果每个$F_i$对该点产生零力矩，则标记为$\pm$；如果至少一个wrench对该点产生正力矩且至少一个wrench对该点产生负力矩，则标记为空白。

这个想法最好通过例子来说明。在图12.21(a)中，基wrench $F_1$通过标记力线左侧的点为$'+'$和线右侧的点为$'-'$来表示。线上的点标记为$\pm$。在图12.21(b)中，添加了另一个基wrench，它可以表示平面摩擦锥的另一条边。只有平面中为两条力线一致标记的点保留其标签；不一致标记的点失去其标签。最后，在图12.21(c)中添加了第三个基wrench。结果是一个标记为$'+'$的单个区域。三个基wrench的非负线性组合可以在平面中创建任何力线，该力线以逆时针方向围绕该区域。不能创建其他wrench。

如果在图12.21(c)中标记为$'+'$的区域周围顺时针添加额外的基wrench，则平面中将没有一致标记的点；四个wrench的正线性张成将是整个wrench空间$\mathbb{R}^3$。

力矩标记表示等价于齐次凸wrench锥表示。图12.21的每个部分(a)、(b)和(c)中的力矩标记区域被正确解释为单个凸区域，就像第12.1.6节的CoR区域一样。

\subsection{力闭合}
\label{subsec:force_closure}

考虑单个可动物体和多个摩擦接触。如果复合wrench锥包含整个wrench空间，使得作用在物体上的任何外部wrench $F_{\text{ext}}$可以由接触力平衡，我们说接触产生力闭合。

我们可以推导出一个简单的力闭合线性测试，对于平面情况是精确的，对于三维情况是近似的。设$F_i$，$i = 1, \ldots, j$，是对应于所有接触的摩擦锥边的wrench。对于平面问题，每个摩擦锥贡献两条边；对于三维问题，每个摩擦锥贡献三条或更多条边，取决于选择的多面体近似（见图12.18(c)）。$n \times j$矩阵$F$的列是$F_i$，其中对于平面问题$n = 3$，对于三维问题$n = 6$。

现在，力闭合的测试与形状闭合的测试相同。如果满足以下条件，则接触产生力闭合：
\begin{itemize}
\item $\text{rank}(F) = n$，且
\item 存在线性规划问题(12.14)的解。
\end{itemize}

在$\mu = 0$的情况下，每个接触只能沿法线方向提供力，力闭合等价于一阶形状闭合。

\subsubsection{力闭合所需的接触数量}
\label{subsubsec:number_of_contacts_force_closure}

对于平面问题，四个接触wrench足以正张成三维wrench空间，这意味着至少两个摩擦接触（每个有两条摩擦锥边）足以实现力闭合。使用力矩标记，我们看到力闭合等价于没有一致的力矩标记。例如，如果两个接触可以通过位于两个摩擦锥内的视线"看到"彼此，则我们有力闭合（图12.22(b)）。

重要的是要注意，力闭合仅仅意味着接触摩擦锥可以产生任何wrench。它不一定意味着物体在存在外部wrench时不会移动。对于图12.22(b)的例子，三角形是否在重力下下落取决于手指之间的内力。如果为手指提供动力的电机无法提供足够的力，或者如果它们被限制为仅在特定方向产生力，则尽管有力闭合，三角形也可能下落。

两个摩擦点接触不足以产生三维物体的力闭合，因为无法产生关于连接两个接触的轴的力矩。然而，可以用至少三个摩擦接触获得力闭合抓取。由[89]给出的一个特别简单且吸引人的结果将三维摩擦抓取的力闭合分析简化为平面力闭合问题。参考图12.23，假设刚体由三个摩擦点接触约束。如果三个接触点恰好共线，那么显然任何关于这条线的力矩都无法由三个接触抵抗。因此我们可以排除这种情况，并假设三个接触点不共线。三个接触然后定义一个唯一的平面$S$，在每个接触点处，出现三种可能性（见图12.24）：
\begin{itemize}
\item 摩擦锥与$S$在平面锥中相交；
\item 摩擦锥与$S$在一条线中相交；
\item 摩擦锥与$S$在一个点处相交。
\end{itemize}

当且仅当每个摩擦锥与$S$在平面锥中相交且$S$也处于平面力闭合时，物体处于力闭合。

\begin{theorem}[12.8]
给定由三个带摩擦的点接触约束的空间刚体，当且仅当每个接触处的摩擦锥与接触的平面$S$在锥中相交且平面$S$处于平面力闭合时，物体处于力闭合。
\end{theorem}

\textbf{证明}。首先，必要性条件——如果空间刚体处于力闭合，那么每个摩擦锥与$S$在平面锥中相交且$S$也处于平面力闭合——很容易验证：如果物体处于空间力闭合，那么$S$（它是物体的一部分）也必须处于平面力闭合。此外，如果甚至一个摩擦锥与$S$在一条线或点中相交，那么将存在无法由抓取抵抗的外部力矩（例如，关于剩余两个接触点之间的线的力矩）。

为了证明充分性条件——如果每个摩擦锥与$S$在平面锥中相交且$S$也处于平面力闭合，那么空间刚体处于力闭合——选择一个固定参考坐标系，使得$S$位于$\hat{x}$–$\hat{y}$平面中，并设$r_i \in \mathbb{R}^3$表示从固定坐标系原点到接触点$i$的向量（见图12.23）。用$f_i \in \mathbb{R}^3$表示接触$i$处的接触力，接触wrench $F_i \in \mathbb{R}^6$的形式为
\begin{equation}
F_i = \begin{bmatrix} m_i \\ f_i \end{bmatrix},
\label{eq:12.17}
\end{equation}
其中每个$m_i = r_i \times f_i$，$i = 1, 2, 3$。用
\begin{equation}
F_{\text{ext}} = \begin{bmatrix} m_{\text{ext}} \\ f_{\text{ext}} \end{bmatrix} \in \mathbb{R}^6
\label{eq:12.18}
\end{equation}
表示任意外部wrench $F_{\text{ext}} \in \mathbb{R}^6$。

力闭合要求存在接触wrench $F_i$，$i = 1, 2, 3$，每个位于其各自的摩擦锥内，使得对于任何外部干扰wrench $F_{\text{ext}}$，满足以下等式：
\begin{equation}
F_1 + F_2 + F_3 + F_{\text{ext}} = 0
\label{eq:12.19}
\end{equation}
或等价地，
\begin{align}
f_1 + f_2 + f_3 + f_{\text{ext}} &= 0, \label{eq:12.20} \\
(r_1 \times f_1) + (r_2 \times f_2) + (r_3 \times f_3) + m_{\text{ext}} &= 0. \label{eq:12.21}
\end{align}

如果每个接触力和力矩，以及外部力和力矩，正交分解为位于由$S$张成的平面（对应于我们选择的参考坐标系中的$\hat{x}$–$\hat{y}$平面）上的分量和其法向子空间$N$（对应于我们选择的参考坐标系中的$\hat{z}$轴）上的分量，则先前的力闭合等式可以写为
\begin{align}
f_{1S} + f_{2S} + f_{3S} &= -f_{\text{ext},S}, \label{eq:12.22} \\
(r_1 \times f_{1S}) + (r_2 \times f_{2S}) + (r_3 \times f_{3S}) &= -m_{\text{ext},S}, \label{eq:12.23} \\
f_{1N} + f_{2N} + f_{3N} &= -f_{\text{ext},N}, \label{eq:12.24} \\
(r_1 \times f_{1N}) + (r_2 \times f_{2N}) + (r_3 \times f_{3N}) &= -m_{\text{ext},N}. \label{eq:12.25}
\end{align}

在接下来的内容中，我们将使用$S$既指对应于$\hat{x}$–$\hat{y}$平面的刚体切片，也指$\hat{x}$–$\hat{y}$平面本身；我们将始终将$N$与$\hat{z}$轴等同。

继续充分性条件的证明，我们现在证明如果$S$处于平面力闭合，那么物体处于空间力闭合。就方程(12.24)和(12.25)而言，我们希望证明，对于任何任意力$f_{\text{ext},S} \in S$、$f_{\text{ext},N} \in N$和任意力矩$m_{\text{ext},S} \in S$、$m_{\text{ext},N} \in N$，存在接触力$f_{iS} \in S$、$f_{iN} \in N$，$i = 1, 2, 3$，满足(12.24)和(12.25)，使得对于每个$i = 1, 2, 3$，接触力$f_i = f_{iS} + f_{iN}$位于摩擦锥$i$内。

首先考虑法向$N$中的力闭合方程(12.24)和(12.25)。给定任意外部力$f_{\text{ext},N} \in N$和外部力矩$m_{\text{ext},S} \in S$，方程(12.24)和(12.25)构成三个未知数的三个线性方程组。根据我们假设三个接触点从不共线，这些方程将始终在$N$中具有唯一解集$\{f_{1N}^*, f_{2N}^*, f_{3N}^*\}$。

由于假设$S$处于平面力闭合，对于任何任意$f_{\text{ext},S} \in S$和$m_{\text{ext},N} \in N$，将存在平面接触力$f_{iS} \in S$，$i = 1, 2, 3$，位于它们各自的平面摩擦锥内，并且也满足方程(12.22)和(12.23)。这个解集不是唯一的：总是可以找到一组内力$\eta_i \in S$，$i = 1, 2, 3$，每个位于其各自的摩擦锥内，满足
\begin{align}
\eta_1 + \eta_2 + \eta_3 &= 0, \label{eq:12.26} \\
(r_1 \times \eta_1) + (r_2 \times \eta_2) + (r_3 \times \eta_3) &= 0. \label{eq:12.27}
\end{align}

（要了解为什么这样的$\eta_i$存在，回想一下，由于假设$S$处于平面力闭合，对于$f_{\text{ext},S} = m_{\text{ext},N} = 0$，方程(12.22)和(12.23)的解必须存在；这些解正是内力$\eta_i$）。注意，这两个方程构成涉及六个变量的三个线性等式约束，因此存在$\{\eta_1, \eta_2, \eta_3\}$的三维线性子空间解。

现在，如果$\{f_{1S}, f_{2S}, f_{3S}\}$满足(12.22)和(12.23)，那么$\{f_{1S} + \eta_1, f_{2S} + \eta_2, f_{3S} + \eta_3\}$也将满足。内力$\{\eta_1, \eta_2, \eta_3\}$可以依次选择为具有足够大的幅度，使得接触力
\begin{align}
f_1 &= f_{1N}^* + f_{1S} + \eta_1, \label{eq:12.28} \\
f_2 &= f_{2N}^* + f_{2S} + \eta_2, \label{eq:12.29} \\
f_3 &= f_{3N}^* + f_{3S} + \eta_3 \label{eq:12.30}
\end{align}
都位于它们各自的摩擦锥内。这完成了充分性条件的证明。

\subsubsection{测量力闭合抓取的质量}
\label{subsubsec:measuring_quality_force_closure}

摩擦力并不总是可重复的。例如，尝试将硬币放在书上并倾斜书。硬币应该在书与水平面成角度$\alpha = \tan^{-1} \mu$时开始滑动。如果你多次进行实验，你可能会发现$\mu$的测量值范围，这是由于难以建模的效应。因此，在选择抓取时，选择最小化实现力闭合所需的摩擦系数的指位置是合理的。

\subsubsection{力和运动自由度的对偶性}
\label{subsubsec:duality_force_motion}

我们对运动学约束和摩擦的讨论应该已经表明，对于任何点接触和接触标签，由该接触引起的物体运动的等式约束数量等于它提供的wrench自由度数量。例如，分离接触$B$对物体运动提供零个等式约束，也不允许接触力。固定接触$R$提供三个运动约束（物体上点的运动被指定）和接触力中的三个自由度：接触wrench锥内部的任何wrench都与接触模式一致。最后，滑动接触$S$提供一个等式运动约束（必须满足物体运动的一个方程以保持接触），并且对于满足约束的给定运动，接触wrench只有一个自由度：摩擦锥边缘上且与滑动方向相反的接触wrench的大小。在平面情况下，$B$、$S$和$R$接触的运动约束和wrench自由度分别为0、1和2。

\section{操作}
\label{sec:manipulation}

到目前为止，我们已经研究了由于一组接触而产生的可行twist和接触力。我们还考虑了两种类型的操作：形状闭合和力闭合抓取。

然而，操作不仅仅包括抓取。它包括几乎任何操作器施加运动或力以实现物体的运动或约束的情况。例子包括在托盘上携带玻璃而不倾倒它们、绕一只脚转动冰箱、沿地板推动沙发、投掷和接球、在振动传送带上运输零件等。赋予机器人超越抓取-携带的操作方法使其能够同时操作多个零件、操作太大而无法抓取或太重而无法提升的零件，甚至通过投掷将零件发送到末端执行器工作空间之外。

为了规划这样的操作任务，我们使用第12.1节的接触运动学约束、第12.2节的库仑摩擦定律以及刚体动力学。将我们自己限制在单个刚体并使用第8章的记号，物体的动力学写为
\begin{equation}
F_{\text{ext}} + \sum_i k_i F_i = G \dot{V} - [\text{ad}_V]^T G V, \quad k_i \geq 0, F_i \in \mathcal{WC}_i,
\label{eq:12.31}
\end{equation}
其中$V$是物体的twist，$G$是其空间惯性矩阵，$F_{\text{ext}}$是由于重力等作用在物体上的外部wrench，$\mathcal{WC}_i$是由于接触$i$作用在物体上的可能wrench集，$\sum_i k_i F_i$是由于接触的wrench。所有wrench都在物体的质心坐标系中表示。现在，给定作用在物体上的一组运动或力控制的接触，以及系统的初始状态，求解物体运动的一种方法如下：
\begin{enumerate}
\item[(a)] 枚举考虑系统当前状态的可能的接触模式集（例如，当前粘附的接触可以转换到滑动或分离）。接触模式由每个接触处的接触标签$R$、$S$和$B$组成。
\item[(b)] 对于每个接触模式，确定是否存在与接触模式和库仑定律一致的接触wrench $\sum_i k_i F_i$，以及与接触模式的运动学约束一致的加速度$\dot{V}$，使得方程(12.31)得到满足。如果是，这个接触模式、接触wrench和物体加速度构成刚体动力学的一致解。
\end{enumerate}

这种"案例分析"可能听起来不寻常；我们不仅仅是在求解一组方程。它还留下了我们可能找到多个一致解，或者可能没有一致解的可能性。事实上，情况就是这样：我们可以定义具有多个解的问题（模糊问题）和没有解的问题（不一致问题）。这种情况有点令人不安；当然任何真实力学问题都恰好有一个解！但这是我们使用完全刚体和库仑摩擦假设所付出的代价。尽管存在零个或多个解的可能性，对于许多问题，上述方法将产生唯一的接触模式和运动。

下面的一些操作任务是准静态的，其中物体的速度和加速度足够小，可以忽略惯性力。接触wrench和外部wrench始终处于力平衡，方程(12.31)简化为
\begin{equation}
F_{\text{ext}} + \sum_i k_i F_i = 0, \quad k_i \geq 0, F_i \in \mathcal{WC}_i.
\label{eq:12.32}
\end{equation}

下面我们用四个例子说明本章的方法。

\textbf{例12.9（由两个手指携带的块）}。考虑重力作用下的平面块，由两个手指支撑，如图12.25(a)所示。一个手指与块之间的摩擦系数为$\mu = 1$，另一个接触是无摩擦的。因此，手指可以施加的wrench锥是$\text{pos}(\{F_1, F_2, F_3\})$，如图12.25(b)中使用力矩标记所示。

我们的第一个问题是静止的手指是否可以使块保持静止。要做到这一点，手指必须提供wrench $F = (m_z, f_x, f_y) = (0, 0, mg)$以平衡由于重力产生的wrench $F_{\text{ext}} = (0, 0, -mg)$，其中$g > 0$。如图12.25(b)所示，这个wrench不在可能的接触wrench的复合锥中。因此，接触模式$RR$不可行，块将相对于手指移动。

现在考虑手指各自向左以$2g$加速的情况。在这种情况下，接触模式$RR$要求块也向左以$2g$加速。产生这个加速度所需的wrench是$(0, -2mg, 0)$。因此，手指必须施加到块的总wrench是$(0, -2mg, 0) - F_{\text{ext}} = (0, -2mg, mg)$。如图12.25(c)、(d)所示，这个wrench位于复合wrench锥内。因此，$RR$（块相对于手指保持静止）是当手指向左以$2g$加速时的解。

这被称为动态抓取——惯性力用于在手指移动时使块压向手指。如果我们计划使用动态抓取操作块，那么为了完整性，我们应该确保除了$RR$之外没有其他接触模式是可行的。

力矩标记对于图形化理解这个问题很方便，但我们也可以代数求解。下手指在$(x, y) = (-3, -1)$处接触块，上手指在$(1, 1)$处接触块。这给出基接触wrench
\begin{align}
F_1 &= \frac{1}{\sqrt{2}}(-4, -1, 1), \\
F_2 &= \frac{1}{\sqrt{2}}(-2, 1, 1), \\
F_3 &= (1, -1, 0).
\end{align}

设手指在$\hat{x}$方向的加速度写为$a_x$。那么，在块相对于手指保持固定（$RR$接触模式）的假设下，方程(12.31)可以写为
\begin{equation}
k_1 F_1 + k_2 F_2 + k_3 F_3 + (0, 0, -mg) = (0, ma_x, 0).
\label{eq:12.33}
\end{equation}

这产生三个未知数$k_1$、$k_2$、$k_3$的三个方程。求解，我们得到
\begin{align}
k_1 &= -\frac{1}{2\sqrt{2}}(a_x + g)m, \\
k_2 &= \frac{1}{2\sqrt{2}}(a_x + 5g)m, \\
k_3 &= -\frac{1}{2}(a_x - 3g)m.
\end{align}

为了使$k_i$非负，我们需要$-5g \leq a_x \leq -g$。对于在此范围内的$\hat{x}$方向手指加速度，动态抓取是一致解。

\textbf{例12.10（米尺技巧）}。尝试这个实验。拿一根米尺（或任何类似的长而光滑的棍子），水平地平衡在你的两个食指上。将你的左手指放在10 cm标记附近，右手指放在60 cm标记附近。质心更靠近你的右手指，但仍然在你的手指之间，因此棍子被支撑。现在，保持你的左手指静止，慢慢地将你的右手指向左移动，直到它们接触。棍子会发生什么？

如果你没有尝试实验，你可能会猜测你的右手指通过棍子的质心下方，此时棍子下落。如果你确实尝试了实验，你看到了不同的情况。让我们看看为什么。

图12.26显示了由两个摩擦手指支撑的棍子。由于所有运动都很慢，我们使用准静态近似，即棍子的加速度为零，因此净接触wrench必须平衡重力wrench。当两个手指一起移动时，棍子必须在一个或两个手指上滑动，以适应手指彼此靠近的事实。图12.26显示了三种不同接触模式的复合wrench锥的力矩标记表示：$RS_r$，其中棍子相对于左手指保持静止，相对于右手指向右滑动；$S_lR$，其中棍子相对于左手指向左滑动，相对于右手指保持静止；$S_lS_r$，其中棍子在两个手指上都滑动。从图中可以清楚地看出，只有$S_lR$接触模式可以提供平衡重力wrench的wrench。换句话说，支撑更多棍子重量的右手指相对于棍子保持固定，而左手指在棍子下滑动。由于右手指在世界坐标系中向左移动，这意味着质心以相同的速度向左移动。这持续到质心在手指之间的一半处，此时棍子转换到$S_lS_r$接触模式，质心保持在手指之间居中，直到它们相遇。棍子永远不会下落。

注意，这个分析依赖于准静态假设。如果你快速移动右手指，很容易使棍子下落；右手指处的摩擦力不足以产生维持粘附接触所需的大棍子加速度。此外，在你的实验中，你可能会注意到，当质心几乎居中时，棍子实际上并没有达到理想的$S_lS_r$接触模式，而是在$S_lR$和$RS_r$接触模式之间快速切换。这是因为静摩擦系数大于动摩擦系数。

\textbf{例12.11（装配的稳定性）}。考虑图12.27中的拱。它在重力下稳定吗？

对于这样的问题，图形平面方法很难使用，因为可能有多个移动物体。相反，我们代数测试所有接触标记为$R$的接触模式的一致性。摩擦锥如图12.27所示。使用这些摩擦锥边的标记，如果存在一致解，即存在$k_i \geq 0$，$i = 1, \ldots, 16$，满足以下九个wrench平衡方程（每个物体三个）：
\begin{align}
\sum_{i=1}^8 k_i F_i + F_{\text{ext}1} &= 0, \\
\sum_{i=9}^{16} k_i F_i + F_{\text{ext}2} &= 0, \\
-\sum_{i=5}^{12} k_i F_i + F_{\text{ext}3} &= 0,
\end{align}
其中$F_{\text{ext}i}$是物体$i$上的重力wrench。最后一组方程来自物体1施加到物体3的wrench等于且相反于物体3施加到物体1的wrench这一事实，对于物体2和3也是如此。

这个线性约束满足问题可以通过多种方法求解，包括线性规划。

\textbf{例12.12（销插入）}。图12.28显示了力控制的平面销在插入过程中与孔的两点接触。还显示了作用在销上的接触摩擦锥和相应的复合wrench锥，使用力矩标记说明。如果力控制器将wrench $F_1$施加到销上，它可能会卡住——孔可能产生平衡$F_1$的接触力。因此，销可能卡在这个位置。然而，如果力控制器将wrench $F_2$施加到销上，接触无法平衡wrench，插入继续进行。

如果两个接触处的摩擦系数足够大，使得两个摩擦锥可以"看到"彼此的基（图12.22(b)），则销处于力闭合，接触可能能够抵抗任何wrench（取决于两个接触之间的内力）。销被称为楔入。

\section{总结}
\label{sec:summary}

\begin{itemize}
\item 解决带摩擦的刚体接触问题需要三个要素：(1)接触运动学，它描述接触刚体的可行运动；(2)接触力模型，它描述可以通过摩擦接触传递的力；以及(3)刚体动力学，如第8章所述。
\item 设两个刚体$A$和$B$在空间坐标系中的$p_A$处点接触。设$\hat{n} \in \mathbb{R}^3$是单位接触法线，指向物体$A$。那么与沿接触法线的单位力相关的空间接触wrench $F$是$F = [([p_A]\hat{n})^T \, \hat{n}^T]^T$。不可穿透约束是
\begin{equation}
F^T (V_A - V_B) \geq 0,
\end{equation}
其中$V_A$和$V_B$是$A$和$B$的空间twist。
\item 粘附或滚动的接触被分配接触标签$R$，滑动的接触被分配标签$S$，正在分离的接触（满足不可穿透约束(12.5)但不满足主动约束(12.6)）被分配标签$B$。对于具有多个接触的物体，接触模式是各个接触处接触标签的连接。
\item 受到多个静止点接触的单个刚体具有满足所有不可穿透约束的齐次（以原点为根部）多面体凸twist锥。
\item $\mathbb{R}^3$中的齐次多面体凸平面twist锥可以等价地表示为平面中带符号旋转中心的凸区域。
\item 如果一组静止接触纯粹通过仅考虑接触法线的运动学分析阻止物体移动，则称物体处于一阶形状闭合。接触$i = 1, \ldots, j$的接触wrench $F_i$正张成$\mathbb{R}^n$，其中$n = 3$在平面情况下，$n = 6$在三维情况下。
\item 平面物体的一阶形状闭合至少需要四个点接触，三维物体的一阶形状闭合至少需要七个点接触。
\item 库仑摩擦定律指出，接触处的切向摩擦力大小$f_t$满足$f_t \leq \mu f_n$，其中$\mu$是摩擦系数，$f_n$是法向力。当接触粘附时，摩擦力可以是满足此约束的任何值。当接触滑动时，$f_t = \mu f_n$，摩擦力的方向与滑动方向相反。
\item 给定作用在物体上的一组摩擦接触，可以通过这些接触传递的wrench是可以通过各个接触传递的wrench的正张成。这些wrench形成齐次凸锥。如果物体是平面的，或者如果物体是三维的但接触摩擦锥由多面体锥近似，则wrench锥也是多面体的。
\item $\mathbb{R}^3$中的齐次凸平面wrench锥可以表示为平面中力矩标记的凸区域。
\item 如果来自静止接触的接触wrench的齐次凸锥是整个wrench空间（$\mathbb{R}^3$或$\mathbb{R}^6$），则物体处于力闭合。如果接触是无摩擦的，力闭合等价于一阶形状闭合。
\end{itemize}

\section{练习}
\label{sec:exercises}

\textbf{练习12.1} 证明不可穿透约束(12.4)等价于约束(12.7)。

\textbf{练习12.2} 将平面twist表示为旋转中心。
\begin{enumerate}
\item[(a)] 考虑两个平面twist $V_1 = (\omega_{z1}, v_{x1}, v_{y1}) = (1, 2, 0)$和$V_2 = (\omega_{z2}, v_{x2}, v_{y2}) = (1, 0, -1)$。在平面坐标系中绘制相应的CoR，并将$\text{pos}(\{V_1, V_2\})$表示为CoR。
\item[(b)] 将$V_1 = (\omega_{z1}, v_{x1}, v_{y1}) = (1, 2, 0)$和$V_2 = (\omega_{z2}, v_{x2}, v_{y2}) = (-1, 0, -1)$的正张成绘制为CoR。
\end{enumerate}

\textbf{练习12.3} 刚体在$p = (1, 2, 3)$处接触，接触法线指向物体$\hat{n} = (0, 1, 0)$。写出由于此接触对物体twist $V$的约束。

\textbf{练习12.4} 在静止约束和物体之间的接触处定义空间坐标系$\{s\}$。接触法线，指向物体，沿着$\{s\}$坐标系的$\hat{z}$轴。
\begin{enumerate}
\item[(a)] 如果接触是无摩擦点接触，写出对物体twist $V$的约束。
\item[(b)] 如果接触是有摩擦点接触，写出对$V$的约束。
\item[(c)] 如果接触是软接触，写出对$V$的约束。
\end{enumerate}

\textbf{练习12.5} 图12.29显示了五个静止"手指"接触物体。物体处于一阶形状闭合，因此力闭合。如果我们拿走一个手指，物体可能仍然处于形状闭合。对于哪些四个手指的子集，物体仍然处于形状闭合？使用图形方法证明你的答案。

\textbf{练习12.6} 当图12.29的三角形仅由手指1接触时，绘制可行twist集作为CoR。用它们的接触标签标记可行的CoR。

\textbf{练习12.7} 当图12.29的三角形仅由手指1和2接触时，绘制可行twist集作为CoR。用它们的接触标签标记可行的CoR。

\textbf{练习12.8} 当图12.29的三角形仅由手指2和3接触时，绘制可行twist集作为CoR。用它们的接触标签标记可行的CoR。

\textbf{练习12.9} 当图12.29的三角形仅由手指1和5接触时，绘制可行twist集作为CoR。用它们的接触标签标记可行的CoR。

\textbf{练习12.10} 当图12.29的三角形仅由手指1、2和3接触时，绘制可行twist集作为CoR。

\textbf{练习12.11} 当图12.29的三角形仅由手指1、2和4接触时，绘制可行twist集作为CoR。

\textbf{练习12.12} 当图12.29的三角形仅由手指1、3和5接触时，绘制可行twist集作为CoR。

\textbf{练习12.13} 再次参考图12.29的三角形。
\begin{enumerate}
\item[(a)] 假设摩擦角$\alpha = 22.5°$（摩擦系数$\mu = 0.41$），使用力矩标记绘制来自接触5的wrench锥。
\item[(b)] 假设摩擦角$\alpha = 45°$（摩擦系数$\mu = 1$），使用力矩标记绘制来自接触5的wrench锥。
\item[(c)] 对于(a)和(b)中的每个摩擦角，确定仅由接触5产生的wrench锥是否包含整个wrench空间。
\end{enumerate}

\textbf{练习12.14} 图12.30显示了由五个无摩擦点接触约束的$4 \times 4$平面正方形。确定哪些接触子集产生形状闭合。

\textbf{练习12.15} 对于图12.30的正方形，假设所有接触都有摩擦系数$\mu = 1$。确定哪些接触子集产生力闭合。

\textbf{练习12.16} 考虑图12.31中显示的平面物体，由三个静止手指接触。每个接触都有摩擦系数$\mu = 1$。使用力矩标记方法确定物体是否处于力闭合。

\textbf{练习12.17} 考虑图12.32中显示的平面物体，由两个静止手指接触。每个接触都有摩擦系数$\mu$。确定实现力闭合所需的最小摩擦系数$\mu$。

\textbf{练习12.18} 考虑图12.33中显示的平面物体，由三个静止手指接触。每个接触都有摩擦系数$\mu = 0.5$。确定物体是否处于力闭合。如果不是，描述物体可以进行的运动。

\textbf{练习12.19} 考虑图12.34中显示的平面物体，由两个静止手指接触。每个接触都有摩擦系数$\mu = 1$。确定物体是否处于力闭合。

\textbf{练习12.20} 考虑图12.35中显示的平面物体，由四个静止手指接触。每个接触都有摩擦系数$\mu = 0.5$。确定物体是否处于力闭合。

\textbf{练习12.21} 考虑图12.36中显示的平面物体，由三个静止手指接触。每个接触都有摩擦系数$\mu = 1$。确定物体是否处于力闭合。如果不是，描述物体可以进行的运动。

\textbf{练习12.22} 考虑图12.37中显示的平面物体，由两个静止手指接触。每个接触都有摩擦系数$\mu = 1$。确定物体是否处于力闭合。如果不是，描述物体可以进行的运动。

\textbf{练习12.23} 考虑图12.38中显示的平面物体，由三个静止手指接触。每个接触都有摩擦系数$\mu = 0.5$。确定物体是否处于力闭合。如果不是，描述物体可以进行的运动。

\textbf{练习12.24} 考虑图12.39中显示的平面物体，由四个静止手指接触。每个接触都有摩擦系数$\mu = 1$。确定物体是否处于力闭合。

\textbf{练习12.25} 考虑图12.40中显示的平面物体，由两个静止手指接触。每个接触都有摩擦系数$\mu = 0.5$。确定物体是否处于力闭合。如果不是，描述物体可以进行的运动。

\textbf{练习12.26} 编程作业。

\textbf{练习12.27} 编程作业。

\textbf{练习12.28} 编程作业。

\textbf{练习12.29} 编程作业。

\textbf{练习12.30} 编程作业。

\textbf{练习12.31} 编程作业。

\end{document}

